\scanpage{110}
\mspage{NA 21}
% archive sheet number: 109
% The typescript is upside down in the supplied scan; reading order below is after rotation by 180 degrees.
% A continuous pen stroke traverses the body of the reused typescript. Its cancellation is represented paragraph by paragraph; nested deletions are local typewritten cancellations.

\deleted{formulation précise en termes d'une notion de ``régularité'' ad hoc (dont Kashiwara ne disposait pas, pas en 1975,76 tout au moins), et enfin dans la démonstration\deleted{\unsure{,}} du théorème, dont les principaux ingrédients avaient \unsure{été} réunis déjà dans sa thèse (dans le cas\unsure{-| }\unsure{soeux|ceux} des \handwritten{${}^{\unsure{\infty|\mathscr{D}}}$}-Modules), à un moment où personne, apparemment, ne se souciait encore le moins du monde de gaspiller un temps précieux à vouloir prouver un énoncé ressenti comme \inserted{accidentel,} marginal.}

\deleted{Le lecteur attentif fera de même les \unsure{ajustements|ajustement} qui s'imposent, quand par la suite il m'arrivera de parler de la ``philosophie de Mebkhout'' des -Modules. Il semble bien, par les documents et témoignages à ma disposition, qu'il ait été le seul, \inserted{entre 1976 et 1980,} à se consacrer corps et âme au développement de cette nouvelle philosophie des coefficients\deleted{\unsure{,}} \deleted{\unsure{mixtes}} en train de naître cahin-caha, et à la faire connaître, et ses publications sont à ce sujet un témoignage éloquent. Cela n'est en aucun cas une raison valable, cependant, pour faire mine d'ignorer pour autant \deleted{\unsure{la}}le \deleted{\unsure{contribution}}\inserted{rôle} de Kashiwara dans l'éclosion de certaines idées importantes, même si ce rôle n'est attesté par aucune publication en bonne et due forme, mais étaient transmises simplement (comme cela a été le cas si souvent pour \unsure{moi-même} dans les années soixante) de bouche à oreille. \deleted{Si donc il \unsure{devait|venait} \illegible{2} chercher un nom à cette philosophie-là (autre que le nom bâtard de ``Riemann-Hilbert'' !), \illegible{4} \unsure{Pour ma part}, j'espère qu'il lui \unsure{restera}.}}

\deleted{Page \unsure{963|953}, deux compléments à la note de haut de page:}

\deleted{($^{x}$) (6 mai 1986) Il y a lieu, visiblement, d'exclure le nom de M. Kashiwara de la liste des participants à la ``mystification-escroquerie'' du Colloque Pervers''. S'il est vrai que ``l'échange de mauvais procédés'' entre Kashiwara et Mebkhout était alors engagé depuis belle lurette, il est vrai que non seulement Kashiwara est étranger à la mise en scène du Colloque Pervers, mais qu'il en est, à part Mebkhout\deleted{\unsure{,}} \deleted{\unsure{le principal}} (et l'ancêtre, il va sans dire), le principal dindon de la farce, vu qu'il n'y est pas plus question de lui, que de l'inconnu de service alias Mebkhout. (Il est vrai que c'était de moindre conséquence pour lui, qui depuis longtemps fait partie de l'establishment mathématique, que pour Mebkhout, pratiquement inconnu et ``seul contre tous''.)}

\deleted{($^{xx}$) (6 mai 1986) Pierre Schapira \inserted{me signale} que l'idée de la transformée de Fourier dans le cadre des -Modules avait été développée (mais pas sous ce nom) dans l'article de \unsure{Seto|Sato},Kawai, Kashiwara de \unsure{1973|1975} déjà cité (note de b. de p. ($\unsure{xxx|***}$) p. 948), donc bien avant les travaux de Malgrange sur ce thème.}
